\documentclass[../DoAn]{subfiles}

\begin{document}

% ============================================================
% MỘT SỐ LƯU Ý VỀ TÀI LIỆU THAM KHẢO
% (chapter* — không đánh số chương)
% ============================================================

Trong quá trình thực hiện đồ án, các tài liệu được tham khảo bao gồm:

\begin{itemize}
    \item \textbf{Văn bản pháp luật Việt Nam} (Luật Khám chữa bệnh, các Thông tư của Bộ Y tế, Quyết định, Chỉ thị của Thủ tướng) được cite trực tiếp từ cổng \textit{vbpl.vn}, \textit{chinhphu.vn} và \textit{thuvienphapluat.vn}.
    \item \textbf{Bài báo khoa học IEEE/ACM} về Blockchain trong Y tế (MedRec MIT 2016, BHEEM 2018, các paper khảo sát blockchain healthcare 2022-2025).
    \item \textbf{Tài liệu kỹ thuật chính thức} của các framework, thư viện được sử dụng (React Native, Tamagui, Solidity, Foundry, The Graph, Web3Auth Sapphire, viem).
    \item \textbf{Tiêu chuẩn quốc tế} (NIST SP 800-38D về AES-GCM, EIP-712 spec của Ethereum Foundation, HL7 FHIR R4, DICOM Standard).
    \item \textbf{Quy chuẩn pháp luật quốc tế tham chiếu} (HIPAA Hoa Kỳ 1996, GDPR Liên minh châu Âu 2016/679) — phục vụ phân tích so sánh.
\end{itemize}

Các trích dẫn trong Quyển tuân thủ phong cách \textbf{IEEE} (\texttt{\textbackslash{}cite} numbered, format trong \texttt{Danh\_sach\_tai\_lieu\_tham\_khao.bib}).

Đối với các đoạn mã nguồn được trích dẫn từ chính dự án, hệ thống cite theo định dạng \texttt{file:line} (ví dụ: \texttt{contracts/src/ConsentLedger.sol:283}) để người đọc dễ truy nguyên gốc trên kho mã nguồn.

\end{document}
